\section{Discussion}
The comparative analysis of HDD, SATA SSD, and NVMe SSD technologies presented 
in this study highlights a clear evolutionary trajectory in storage system design, 
driven primarily by the need to eliminate mechanical bottlenecks and better exploit 
the parallelism available in modern processors. While HDDs remain relevant for high-capacity, 
cost-sensitive storage due to their low cost per gigabyte, their reliance on mechanical 
components fundamentally limits their performance and reliability compared to flash-based alternatives.

The transition from SATA-based SSDs to NVMe represents more than an incremental interface 
upgrade; it reflects a re-architecture of how the operating system communicates with storage. 
By moving from a single command queue with a shallow queue depth to tens of thousands of 
parallel queues, NVMe allows the operating system to fully utilize the throughput offered 
by the PCIe bus and the multi-core nature of modern CPUs. This has direct implications for 
I/O scheduling, since traditional seek-time-minimizing algorithms designed for HDDs become 
largely irrelevant for both SATA SSDs and NVMe drives, prompting operating systems to adopt 
simpler, fairness-oriented scheduling policies.

At the same time, the study shows that all three technologies still depend on the operating 
system for essential services such as partitioning, file system management, and error 
handling, even though the specific mechanisms differ significantly. For flash-based storage, 
TRIM support, wear leveling, and garbage collection place additional responsibilities on 
both the drive controller and the operating system to preserve long-term performance and 
endurance, responsibilities that have no HDD equivalent.

From a practical standpoint, the findings suggest that storage technology selection should 
be guided by workload characteristics rather than by a single performance metric. High-capacity 
archival systems continue to benefit from HDDs, latency-sensitive consumer and enterprise 
workloads increasingly favor SATA or NVMe SSDs, and data-intensive applications such as 
databases, virtualization, and artificial intelligence workloads benefit most from NVMe's 
high queue depth and low latency. This reinforces the broader trend in the industry toward 
hybrid storage architectures that combine multiple technologies to balance cost, capacity, 
and performance.

